\documentclass[10pt,a4paper]{amsart}
\usepackage{amsmath,amssymb,amsthm,graphicx}
% Retained minimal preamble for the Stacks Project chapter more-algebra.tex (More on Algebra).
% The chapter file pinned-original-source.tex carries only \input{preamble} on its first line, so no
% preamble text lives inside the chapter itself.  The Stacks Project uses one single shared preamble for
% every chapter; the snapshot stages it identically as sources/stacks-chow/preamble.tex and
% sources/stacks-etale-cohomology/preamble.tex.  The block below is copied mechanically, line for line,
% from sources/stacks-chow/preamble.tex lines 14-18 (the Xy-pic setup) and 154-217 (the native theorem
% declarations together with the mathematical macros).  Nothing was added, removed or redefined; the
% wrapper class and the amsmath/amssymb/amsthm/graphicx base packages are supplied by the assembler.
% Chapter: more-algebra.tex -- More on Algebra.
% For commutative diagrams we use Xy-pic
\usepackage[all]{xy}

% We use 2cell for 2-commutative diagrams.
\xyoption{2cell}

% Theorem environments.
%
\theoremstyle{plain}
\newtheorem{theorem}[subsection]{Theorem}
\newtheorem{proposition}[subsection]{Proposition}
\newtheorem{lemma}[subsection]{Lemma}

\theoremstyle{definition}
\newtheorem{definition}[subsection]{Definition}
\newtheorem{example}[subsection]{Example}
\newtheorem{exercise}[subsection]{Exercise}
\newtheorem{situation}[subsection]{Situation}

\theoremstyle{remark}
\newtheorem{remark}[subsection]{Remark}
\newtheorem{remarks}[subsection]{Remarks}

\numberwithin{equation}{subsection}

% Macros
%
\def\lim{\mathop{\mathrm{lim}}\nolimits}
\def\colim{\mathop{\mathrm{colim}}\nolimits}
\def\Spec{\mathop{\mathrm{Spec}}}
\def\Hom{\mathop{\mathrm{Hom}}\nolimits}
\def\Ext{\mathop{\mathrm{Ext}}\nolimits}
\def\SheafHom{\mathop{\mathcal{H}\!\mathit{om}}\nolimits}
\def\SheafExt{\mathop{\mathcal{E}\!\mathit{xt}}\nolimits}
\def\Sch{\mathit{Sch}}
\def\Mor{\mathop{\mathrm{Mor}}\nolimits}
\def\Ob{\mathop{\mathrm{Ob}}\nolimits}
\def\Sh{\mathop{\mathit{Sh}}\nolimits}
\def\NL{\mathop{N\!L}\nolimits}
\def\CH{\mathop{\mathrm{CH}}\nolimits}
\def\proetale{{pro\text{-}\acute{e}tale}}
\def\etale{{\acute{e}tale}}
\def\QCoh{\mathit{QCoh}}
\def\Ker{\mathop{\mathrm{Ker}}}
\def\Im{\mathop{\mathrm{Im}}}
\def\Coker{\mathop{\mathrm{Coker}}}
\def\Coim{\mathop{\mathrm{Coim}}}

% Boxtimes
%
\DeclareMathSymbol{\boxtimes}{\mathbin}{AMSa}{"02}

%
% Macros for moduli stacks/spaces
%
\def\QCohstack{\mathcal{QC}\!\mathit{oh}}
\def\Cohstack{\mathcal{C}\!\mathit{oh}}
\def\Spacesstack{\mathcal{S}\!\mathit{paces}}
\def\Quotfunctor{\mathrm{Quot}}
\def\Hilbfunctor{\mathrm{Hilb}}
\def\Curvesstack{\mathcal{C}\!\mathit{urves}}
\def\Polarizedstack{\mathcal{P}\!\mathit{olarized}}
\def\Complexesstack{\mathcal{C}\!\mathit{omplexes}}
% \Pic is the operator that assigns to X its picard group, usage \Pic(X)
% \Picardstack_{X/B} denotes the Picard stack of X over B
% \Picardfunctor_{X/B} denotes the Picard functor of X over B
\def\Pic{\mathop{\mathrm{Pic}}\nolimits}
\def\Picardstack{\mathcal{P}\!\mathit{ic}}
\def\Picardfunctor{\mathrm{Pic}}
\def\Deformationcategory{\mathcal{D}\!\mathit{ef}}
\usepackage{cleveref}
\crefname{originallemma}{original lemma}{original lemmas}
\Crefname{originallemma}{Original lemma}{Original lemmas}
\makeatletter
\newlabel{algebra-lemma-cover}{{00EO}{source}{}{sourcepointer-1}{}}
\newlabel{algebra-lemma-cover@cref}{{[originallemma][0][]00EO}{[1][0][]source}{}{}{}}
\makeatother
\title{Splitting a ring over the open pieces of a finitely generated ideal by adding one surjective component}
\author{The Stacks Project Authors}
\date{stacks-project:more-algebra.tex:local-snapshot-20261012; GNU FDL 1.2 or later}
\begin{document}
\maketitle
\setcounter{subsection}{1}
\noindent\textit{Editorial note.} All delivered mathematical text is the author's own, copied from the pinned Stacks Project chapter file more-algebra.tex; no statement or proof step was rewritten, reordered or removed. The counted statement is the labeled lemma environment at source lines 31983-31990 (canonical tag 0C2A) and the counted proof is the proof environment at source lines 31992-32027, which immediately follows it and is its own proof.\par\medskip
\noindent\textit{Editorial note.} The proof refers once, in its last paragraph, to the labelled lemma algebra-lemma-cover of another chapter for the glueing of functions. That label is not part of this chapter file, so the card carries one navigation stub whose pointer is the canonical tag list line of the snapshot; the stub only resolves the cross reference in the compiled document and carries no mathematics.\par\medskip
\noindent\textit{Editorial note.} The lead-in sentence at source line 31981, calling the statement a simple glueing lemma, is outside the counted spans, as is the remark at source lines 32029-32030 pointing to another lemma for the construction of quotients by idempotents.\par\medskip
\noindent\textit{Editorial note.} The chapter file pinned-original-source.tex carries only an input of the shared preamble on its first line, so no preamble text lives inside the chapter. The retained file sources/stacks-more-algebra/retained-preamble.tex is a mechanical, line-for-line copy of the project's single shared preamble as staged in the snapshot (Xy-pic setup plus the native theorem declarations and mathematical macros). Nothing was added, removed or redefined, and the excerpt serves only as preamble, never as a span coordinate.\par\medskip
\noindent\textit{Editorial note.} The source URL is the source-level Stacks Project URL, as required when more than one card is produced from one source; the per-claim tag is carried by claim\_key and by this note. The counted statement and proof cite no bibliography entry, so this card needs no bibliography.\par\medskip
\noindent\textit{Allowed standard prerequisites (editorial).}\begin{itemize}
\item Localization at an idempotent splits a ring as a product of the localization at the idempotent and the quotient by the ideal generated by it. used at proof lines 32015-32018, where C'\_\{f\_i\} = A\_\{f\_i\}/(e\_i) and the decomposition B = B\_e times B/(e) for an idempotent e is applied. Identifies the localization of the constructed component C' at each generator with the complementary factor of C\_\{f\_i\}, which is what produces the product C times C'.
\item An idempotent of a localization A\_f can be written in the form a/f\textasciicircum{}n, and the representative can be adjusted so that it is itself an idempotent after multiplying by a power of f. used at proof lines 31997-32001, where e\_i = a\_i/f\_i\textasciicircum{}n is written and the relation a\_i\textasciicircum{}2 = f\_i\textasciicircum{}n a\_i is arranged. Produces the elements a\_i in A whose ideal defines C' and makes the comparison of the two components at f\_i computable.
\item If two elements agree after inverting f\_i f\_j for all i and j then they agree after inverting the product of the generators, and an isomorphism that holds after localizing at f\_1, ..., f\_r holds after localizing at any element of the ideal they generate. used at proof lines 32002-32006 and 32020-32026. Gives the bounded relations (f\_i f\_j a\_j)\textasciicircum{}N (f\_j\textasciicircum{}n a\_i - f\_i\textasciicircum{}n a\_j) = 0 and, at the end, transfers the isomorphism from the generators to every element of I.
\end{itemize}\medskip
\ifdefined\hypertarget\hypertarget{sourcepointer-1}{}\fi
\noindent\textit{Source reference.} Original lemma 00EO: sources/stacks-chow/tags.txt, source lines 542--542. The final sentence uses the glueing of local isomorphisms of another chapter to pass from the generators f\_1, ..., f\_r to every element f of I; the whole construction of C', the exponent bookkeeping and the verification at each generator are inside the counted span, so the reference is a pointer to a standard glueing statement and not a substitute for the argument.\par\medskip
\begin{lemma}
\label{lemma-glueing-sum-components-open}
Let $A$ be a ring and let $I$ be a finitely generated ideal.
Let $A \to C$ be a ring map such that for all $f \in I$
the ring map $A_f \to C_f$ is localization at an idempotent.
Then there exists a surjection $A \to C'$ such that
$A_f \to (C \times C')_f$ is an isomorphism for all $f \in I$.
\end{lemma}

\begin{proof}
Choose generators $f_1, \ldots, f_r$ of $I$. Write
$$
C_{f_i} = (A_{f_i})_{e_i}
$$
for some idempotent $e_i \in A_{f_i}$. Write $e_i = a_i/f_i^n$
for some $a_i \in A$ and $n \geq 0$; we may use the same $n$
for all $i = 1, \ldots, r$. After replacing $a_i$ by $f_i^ma_i$
and $n$ by $n + m$ for a suitable $m \gg 0$, we may assume
$a_i^2 = f_i^n a_i$ for all $i$. Since $e_i$ maps to $1$ in
$C_{f_if_j} = (A_{f_if_j})_{e_j} = A_{f_if_ja_j}$ we see that
$$
(f_if_ja_j)^N(f_j^n a_i  - f_i^na_j) = 0
$$
for some $N$ (we can pick the same $N$ for all pairs $i, j$). Using
$a_j^2 = f_j^na_j$ this gives
$$
f_i^{N + n} f_j^{N + nN} a_j = f_i^N f_j^{N + n} a_ia_j^N
$$
After increasing $n$ to $n + N + nN$ and replacing $a_i$ by $f_i^{N + nN}a_i$
we see that $f_i^n a_j$ is in the ideal of $a_i$ for all pairs $i, j$.
Let $C' = A/(a_1, \ldots, a_r)$. Then
$$
C'_{f_i} = A_{f_i}/(a_i) = A_{f_i}/(e_i)
$$
because $a_j$ is in the ideal generated by $a_i$ after inverting $f_i$.
Since for an idempotent $e$ of a ring $B$ we have $B = B_e \times B/(e)$
we see that the conclusion of the lemma holds for $f$ equal to one
of $f_1, \ldots, f_r$. Using glueing of functions, in the form of
Algebra, Lemma \ref{algebra-lemma-cover},
we conclude that the result holds for all $f \in I$.
Namely, for $f \in I$ the elements $f_1, \ldots, f_r$
generate the unit ideal in $A_f$ so $A_f \to (C \times C')_f$
is an isomorphism if and only if this is the case after localizing
at $f_1, \ldots, f_r$.
\end{proof}

\end{document}
